% GENERATED FILE. Edit canonical lessons, then run npm run generate:books.
% canonical-source-hash: fnv1a64:b429856673f32f38
% canonical-lessons: MR-C337-dantkatha, MR-C337-povada, MR-C337-amar, MR-C337-prakhyat, MR-C337-kavi

\chapter{Legends and Ballads}
\label{ch:mr-legends-and-ballads}
\hlchaptermodality{337}

\begin{chapteropening}
\textbf{By the end of this chapter:} \emph{I can talk about legends, ballads and famous poets.}

Complete the last lesson of chapter 337: I can talk about legends, ballads and famous poets.
\end{chapteropening}

\section[\texorpdfstring{\mr{दंतकथा} (dantkathā)}{दंतकथा (dantkathā)}]{\mr{दंतकथा} (dantkathā) — a legend}

\label{lesson:MR-C337-dantkatha}



\begin{quote}
Before the new one: say the Marathi for an underground passage, then the Marathi for a treasure.
\end{quote}

\subsection*{You'll want to know: \mr{दंतकथा}}
\textbf{\mr{दंतकथा}} — \emph{dantkathā} — \textquotedblleft{}a legend\textquotedblright{}.

\begin{quote}
\mr{या} \mr{तलावाबद्दल} \mr{एक} \mr{जुनी} \mr{दंतकथा} \mr{आहे}.
\end{quote}

\begin{quote}
\emph{yā talāvābaddal ek junī dantkathā āhe.}
\end{quote}

\begin{quote}
There is an old legend about this lake.
\end{quote}

Feminine.

A tale passed on by word of mouth.

\begin{cousinweb}
\textbf{\mr{दंतकथा}} (\emph{dantkathā}) is built from Sanskrit \emph{danta}, 'a tooth', and \emph{kathā}, 'a story' (Monier-Williams); it is usually explained as a tale passed from mouth to mouth.
\end{cousinweb}

\begin{grammarlens}[title={What you've built}]
One more word for telling what happened.
\end{grammarlens}

\subsection*{Your turn}
\begin{itemize}
\raggedright
  \item \emph{Say it:} \emph{dantkathā}
  \item \emph{Say it:} \emph{dantkathā}, once more
  \item \emph{Say it:} \emph{yā talāvābaddal ek junī dantkathā āhe.} — There is an old legend about this lake.
  \item \emph{Recall:} say the Marathi for an underground passage, then the Marathi for a treasure, then say \emph{dantkathā} again
\end{itemize}

\begin{tcolorbox}[breakable,colback=teal!4,colframe=teal!35!black,title={Before you move on}]
What does \mr{दंतकथा} mean? (\textquotedblleft{}A legend\textquotedblright{}.) Say it once more.
\end{tcolorbox}
\section[\texorpdfstring{\mr{पोवाडा} (povāḍā)}{पोवाडा (povāḍā)}]{\mr{पोवाडा} (povāḍā) — a heroic ballad}

\label{lesson:MR-C337-povada}



\begin{quote}
Before the new one: say the Marathi for a treasure, then the Marathi for a legend.
\end{quote}

\subsection*{You'll want to know: \mr{पोवाडा}}
\textbf{\mr{पोवाडा}} — \emph{povāḍā} — \textquotedblleft{}a heroic ballad\textquotedblright{}.

\begin{quote}
\mr{जत्रेत} \mr{एक} \mr{गायक} \mr{शूर} \mr{सरदाराचा} \mr{पोवाडा} \mr{गात} \mr{होता}.
\end{quote}

\begin{quote}
\emph{jatret ek gāyak śūr sardārācā povāḍā gāt hotā.}
\end{quote}

\begin{quote}
At the fair a singer was singing a ballad of a brave chieftain.
\end{quote}

Masculine.

A Marathi ballad of heroic deeds.

\begin{cousinweb}
\textbf{\mr{पोवाडा}} (\emph{povāḍā}) is a tadbhava: CDIAL lists Marathi \emph{povāḍā}, 'a historical ballad', under Sanskrit \emph{pravāda}, 'a speaking forth; talk, report' (CDIAL 8795).
\end{cousinweb}

\begin{grammarlens}[title={What you've built}]
One more word for telling what happened.
\end{grammarlens}

\subsection*{Your turn}
\begin{itemize}
\raggedright
  \item \emph{Say it:} \emph{povāḍā}
  \item \emph{Say it:} \emph{povāḍā}, once more
  \item \emph{Say it:} \emph{jatret ek gāyak śūr sardārācā povāḍā gāt hotā.} — At the fair a singer was singing a ballad of a brave chieftain.
  \item \emph{Recall:} say the Marathi for a treasure, then the Marathi for a legend, then say \emph{povāḍā} again
\end{itemize}

\begin{tcolorbox}[breakable,colback=teal!4,colframe=teal!35!black,title={Before you move on}]
What does \mr{पोवाडा} mean? (\textquotedblleft{}A heroic ballad\textquotedblright{}.) Say it once more.
\end{tcolorbox}
\section[\texorpdfstring{\mr{अमर} (amar)}{अमर (amar)}]{\mr{अमर} (amar) — immortal; undying}

\label{lesson:MR-C337-amar}



\begin{quote}
Before the new one: say the Marathi for a legend, then the Marathi for a heroic ballad.
\end{quote}

\subsection*{You'll want to know: \mr{अमर}}
\textbf{\mr{अमर}} — \emph{amar} — \textquotedblleft{}immortal; undying\textquotedblright{}.

\begin{quote}
\mr{त्या} \mr{वीराचं} \mr{नाव} \mr{इतिहासात} \mr{अमर} \mr{झालं}.
\end{quote}

\begin{quote}
\emph{tyā vīrācaṁ nāv itihāsāt amar jhālaṁ.}
\end{quote}

\begin{quote}
That hero's name became immortal in history.
\end{quote}

\mr{अमर} keeps one form: \mr{अमर} \mr{कथा}, \mr{अमर} \mr{नाव}.

\begin{cousinweb}
\textbf{\mr{अमर}} (\emph{amar}) is a tatsama, taken from Sanskrit \emph{amara}, 'undying, immortal' (Monier-Williams).
\end{cousinweb}

\begin{grammarlens}[title={What you've built}]
One more word for telling what happened.
\end{grammarlens}

\subsection*{Your turn}
\begin{itemize}
\raggedright
  \item \emph{Say it:} \emph{amar}
  \item \emph{Say it:} \emph{amar}, once more
  \item \emph{Say it:} \emph{tyā vīrācaṁ nāv itihāsāt amar jhālaṁ.} — That hero's name became immortal in history.
  \item \emph{Recall:} say the Marathi for a legend, then the Marathi for a heroic ballad, then say \emph{amar} again
\end{itemize}

\begin{tcolorbox}[breakable,colback=teal!4,colframe=teal!35!black,title={Before you move on}]
What does \mr{अमर} mean? (\textquotedblleft{}Immortal; undying\textquotedblright{}.) Say it once more.
\end{tcolorbox}
\section[\texorpdfstring{\mr{प्रख्यात} (prakhyāt)}{प्रख्यात (prakhyāt)}]{\mr{प्रख्यात} (prakhyāt) — renowned}

\label{lesson:MR-C337-prakhyat}



\begin{quote}
Before the new one: say the Marathi for a heroic ballad, then the Marathi for immortal.
\end{quote}

\subsection*{You'll want to know: \mr{प्रख्यात}}
\textbf{\mr{प्रख्यात}} — \emph{prakhyāt} — \textquotedblleft{}renowned\textquotedblright{}.

\begin{quote}
\mr{ते} \mr{एक} \mr{प्रख्यात} \mr{इतिहासकार} \mr{आहेत}.
\end{quote}

\begin{quote}
\emph{te ek prakhyāt itihāskār āhet.}
\end{quote}

\begin{quote}
He is a renowned historian.
\end{quote}

\mr{प्रख्यात} keeps one form; \mr{प्रसिद्ध} is the everyday word.

\begin{cousinweb}
\textbf{\mr{प्रख्यात}} (\emph{prakhyāt}) is a tatsama, taken from Sanskrit \emph{prakhyāta}, 'known, celebrated' (Monier-Williams).
\end{cousinweb}

\begin{grammarlens}[title={What you've built}]
One more word for telling what happened.
\end{grammarlens}

\subsection*{Your turn}
\begin{itemize}
\raggedright
  \item \emph{Say it:} \emph{prakhyāt}
  \item \emph{Say it:} \emph{prakhyāt}, once more
  \item \emph{Say it:} \emph{te ek prakhyāt itihāskār āhet.} — He is a renowned historian.
  \item \emph{Recall:} say the Marathi for a heroic ballad, then the Marathi for immortal, then say \emph{prakhyāt} again
\end{itemize}

\begin{tcolorbox}[breakable,colback=teal!4,colframe=teal!35!black,title={Before you move on}]
What does \mr{प्रख्यात} mean? (\textquotedblleft{}Renowned\textquotedblright{}.) Say it once more.
\end{tcolorbox}
\section[\texorpdfstring{\mr{कवी} (kavī)}{कवी (kavī)}]{\mr{कवी} (kavī) — a poet}

\label{lesson:MR-C337-kavi}



\begin{quote}
Before the new one: say the Marathi for immortal, then the Marathi for renowned.
\end{quote}

\subsection*{You'll want to know: \mr{कवी}}
\textbf{\mr{कवी}} — \emph{kavī} — \textquotedblleft{}a poet\textquotedblright{}.

\begin{quote}
\mr{एका} \mr{कवीने} \mr{आमच्या} \mr{गावावर} \mr{सुंदर} \mr{कविता} \mr{लिहिली}.
\end{quote}

\begin{quote}
\emph{ekā kavīne āmacyā gāvāvar sundar kavitā lihilī.}
\end{quote}

\begin{quote}
A poet wrote a beautiful poem about our village.
\end{quote}

Masculine.

Plural \mr{कवी}.

\begin{cousinweb}
\textbf{\mr{कवी}} (\emph{kavī}) is a tatsama, taken from Sanskrit \emph{kavi}, 'a wise man, a poet' (Monier-Williams).
\end{cousinweb}

\begin{grammarlens}[title={What you've built}]
One more word for telling what happened.
\end{grammarlens}

\subsection*{Your turn}
\begin{itemize}
\raggedright
  \item \emph{Say it:} \emph{kavī}
  \item \emph{Say it:} \emph{kavī}, once more
  \item \emph{Say it:} \emph{ekā kavīne āmacyā gāvāvar sundar kavitā lihilī.} — A poet wrote a beautiful poem about our village.
  \item \emph{Recall:} say the Marathi for immortal, then the Marathi for renowned, then say \emph{kavī} again
\end{itemize}

\begin{tcolorbox}[breakable,colback=teal!4,colframe=teal!35!black,title={Before you move on}]
What does \mr{कवी} mean? (\textquotedblleft{}A poet\textquotedblright{}.) Say it once more.
\end{tcolorbox}
